\documentclass[11pt]{article}
\usepackage[margin=1in]{geometry}
\usepackage{amsmath,amssymb,amsfonts}
\usepackage{graphicx}
\usepackage{booktabs}
\usepackage{hyperref}

\title{Closed-Form ZK: Efficient Temporal Proofs\\via CfC Neural Networks}

\author{
Tristan Stoltz\\
Luminous Dynamics\\
\texttt{tristan.stoltz@evolvingresonantcocreationism.com}
}

\begin{document}
\maketitle

\begin{abstract}
We show that Closed-form Continuous-time (CfC) neural networks enable significantly more efficient zero-knowledge temporal proofs than ODE-based alternatives. The CfC closed-form solution $h(t) = (1-\sigma(f)) \cdot h(0) + \sigma(f) \cdot x_\infty$ decomposes into exactly \textbf{1 AND constraint per neuron per timestep} in Binius (a binary-field STARK system), with all additive state updates free. With in-circuit sigmoid approximation, this increases to 2 AND per step. By contrast, Liquid Time-Constant (LTC) networks with RK4 integration require $\sim$16 AND per step. We benchmark real Binius proofs at two scales: 8 neurons $\times$ 10 timesteps (96 AND, 493~ms) and 64 neurons $\times$ 100 timesteps (7,360 AND, 954~ms without sigmoid; 13,760 AND, 680~ms with sigmoid). These results demonstrate that CfC's algebraic structure, combined with binary-field proof systems, makes verifiable temporal neural computation practical for the first time.
\end{abstract}

\section{Introduction}

Temporal neural networks model continuous-time dynamics in time-series data. Liquid Time-Constant (LTC) networks~\cite{hasani2021ltc} and their closed-form counterpart CfC~\cite{hasani2022cfc} have shown strong performance in autonomous driving, medical monitoring, and robotics. A growing need exists for \emph{verifiable} temporal computation: proving that a neural network processed a time series correctly without revealing the input data.

Zero-knowledge proofs (ZKPs) can provide this guarantee, but the computational cost of proving ODE integration is prohibitive. Each RK4 substep requires multiple multiplications per neuron, and $S \approx 16$ substeps per timestep is typical for numerical accuracy. For $N$ neurons over $T$ timesteps, an ODE-based LTC requires $O(16 \cdot N \cdot T)$ non-linear constraints.

CfC networks eliminate this overhead with a closed-form solution that replaces iterative integration with a single expression per timestep. We show that this algebraic simplification has a direct correspondence in binary-field STARK proof systems (Binius~\cite{diamond2023binius}), where the CfC update decomposes into:
\begin{align}
\delta &= x_\infty \oplus h \quad &\text{(XOR: free in Binius)} \\
p &= \sigma \wedge \delta \quad &\text{(AND: 1 constraint)} \\
h' &= h \oplus p \quad &\text{(XOR: free)}
\end{align}

\section{Background}

\textbf{CfC Networks.} The CfC solution~\cite{hasani2022cfc} is:
$$h(t) = (1 - \sigma(f(t))) \cdot h(0) + \sigma(f(t)) \cdot x_\infty$$
where $\sigma$ is sigmoid and $f(t) = -\Delta t / \tau$. This is an $O(1)$ expression per neuron per timestep --- no numerical integration needed.

\textbf{Binius.} A STARK proof system over binary tower fields $\mathrm{GF}(2^k)$~\cite{diamond2023binius}. Addition (XOR) is native field addition and costs zero non-linear constraints. Only AND/MUL operations incur constraints.

\section{CfC Temporal Proof Circuit}

\subsection{Without Sigmoid (Witness-Supplied)}

When $\sigma$ is provided as a precomputed witness value, the CfC update requires exactly \textbf{1 AND per neuron per timestep}: the multiplication $\sigma \wedge \delta$. Both XOR operations are free in Binius.

\subsection{With In-Circuit Sigmoid}

We implement sigmoid as a piecewise mask approximation:
$$\sigma(f) \approx f \wedge \texttt{0xFF}$$
This extracts the low 8 bits of $f$ as a crude sigmoid output, requiring 1 additional AND per step. Total: \textbf{2 AND per neuron per timestep}.

A production-grade piecewise linear sigmoid (5--10 segments) would require 5--10 AND per step, still far fewer than the $\sim$16 for ODE-based LTC.

\section{Experimental Results}

All benchmarks use Binius64~\cite{binius64} with real cryptographic proof generation and verification. Hardware: AMD Ryzen, 32~GB RAM, NixOS.

\begin{table}[h]
\centering
\caption{CfC Temporal Proofs in Binius (Real Proofs)}
\label{tab:cfc}
\begin{tabular}{@{}lrrrr@{}}
\toprule
\textbf{Config} & \textbf{AND} & \textbf{XOR (free)} & \textbf{Prove} & \textbf{Size} \\
\midrule
\multicolumn{5}{l}{\emph{Without sigmoid (witness-supplied $\sigma$)}} \\
8N $\times$ 10T & 96 & 160 & 493 ms & 478 KB \\
64N $\times$ 100T & 7,360 & 12,800 & 954 ms & 156 KB \\
\midrule
\multicolumn{5}{l}{\emph{With in-circuit sigmoid ($\sigma = f \wedge$ \texttt{0xFF})}} \\
64N $\times$ 100T & 13,760 & 12,800 & 680 ms & 166 KB \\
\midrule
\multicolumn{5}{l}{\emph{Estimated: LTC + RK4 (16 substeps per timestep)}} \\
64N $\times$ 100T & $\sim$117,760 & --- & $\sim$10s+ & --- \\
\bottomrule
\end{tabular}
\end{table}

The 64-neuron, 100-timestep CfC proof generates in under 1 second with 7,360 AND constraints. An equivalent LTC+RK4 circuit would require $\sim$117,760 AND constraints ($16\times$ more).

\section{Limitations}

\begin{itemize}
    \item Our sigmoid approximation ($f \wedge \texttt{0xFF}$) is crude. A numerically accurate piecewise linear sigmoid would add 5--10 AND per step instead of 1.
    \item We benchmark temporal evolution in isolation. A full neural inference proof would also need to prove input encoding and output decoding.
    \item The LTC+RK4 comparison is estimated, not measured. Implementing a full ODE-based temporal proof circuit is future work.
    \item Binius64 is pre-1.0 software; API and performance may change.
\end{itemize}

\section{Related Work}

zkRNN~\cite{zkrnn2026} proves RNN execution but does not exploit closed-form temporal solutions. SUMMER~\cite{summer2025} adds recursive ZKPs for RNN training. Neither addresses the CfC/LTC distinction or binary-field alignment. Our companion paper~\cite{binius_hdc} establishes the broader algebraic alignment between HDC and binary-field STARKs.

\section{Conclusion}

CfC's closed-form solution translates directly into efficient binary-field STARK circuits: 1--2 AND per neuron per timestep, with all additive operations free. This is $\sim$16$\times$ fewer constraints than ODE-based temporal networks, making verifiable neural computation practical for temporal data processing.

\bibliographystyle{plain}
\begin{thebibliography}{9}

\bibitem{hasani2022cfc}
R.~Hasani et~al., ``Closed-form continuous-time neural networks,'' \emph{Nature Machine Intelligence}, vol.~4, pp.~992--1003, 2022.

\bibitem{hasani2021ltc}
R.~Hasani et~al., ``Liquid time-constant networks,'' in \emph{Proc. AAAI}, 2021.

\bibitem{diamond2023binius}
B.~Diamond and J.~Posen, ``Succinct arguments over towers of binary fields,'' \emph{ePrint} 2023/1784, 2023.

\bibitem{binius64}
Irreducible, Inc., ``Binius64,'' \url{https://github.com/IrreducibleOSS/binius64}, 2025.

\bibitem{zkrnn2026}
M.~Zarinjouei et~al., ``zkRNN,'' \emph{ePrint} 2026/073, 2026.

\bibitem{summer2025}
``SUMMER: Recursive ZKPs for RNN training,'' \emph{ePrint} 2025/1688, 2025.

\bibitem{binius_hdc}
T.~Stoltz, ``Binary-field STARKs for hyperdimensional computing,'' 2026. Companion paper.

\end{thebibliography}

\end{document}
